\section{有限計算の証明依存、現行の再現可能性と未解決箇所}
\label{sec:finite-calculation-scope}

本節は本文の数学的結果に必要な計算上の入力と検証範囲をまとめる。
本文の幾何学的、圏論的証明はそれぞれの命題・補題で与え、
以下の整数計算・有限標本の検査から一般定理が自動的に従うとは主張しない。

\subsection{構造的入力と整数計算の区別}
指標・共変量の計算にはEuler積、Newton--Adams再帰、Schur--Weyl分解、
Pieri則およびKoszul複体を用いる。
接触方向の線形独立性の主張には、有理数体または有限体上の
\emph{非零小行列式}による証明書と、
その行列の構成を再現するプログラムの両方が必要である。
有限体での非零小行列式は整数行列の特性零での下界を与えるが、
上界や普遍的同一性を代用しない。
Tate再帰の新規計算と既存JSON値の照合、完全点の値と形式jetの値を混同しない。

過去の研究で記録された具体的な計算規模は、
18次の通常接触94方向と追加60方向、
20次付近の164接触方向と359 Koszul方向、
22次の429独立方向、
24次の1091接触方向（1190281成分の選択小行列）
などである。これらは、該当する本文の数学的証明の代わりに
\emph{今回再実行した検算結果}として引用する数字ではない。
元の完全な整数入力と実行スクリプトが現行Git treeに揃っていないものは
独立検証を未実施とし、数値を証明済みの一般的階数と同一視しない。

\subsection{現行本文で保持する定理と未決定事項}
全次数normal成分の消滅、横断射影の頂点支持、
normal消滅の仮定の下での明示的有限決定性は
\ref{sec:finite-calculation-scope}より前の数学的本文に証明が記載されている。
ただし当該証明の独立監査は完了していない。
接触正規形が26次開始になるという原稿の主張を、
未検証の26次係数から28次開始へ強化しない。

$E_{\mathrm{vert}}=t^{10}$、$\mathcal K_{26}^{\mathrm{nor}}=0$と
十次正規形への同変代数化は、定理\ref{thm:pf-full-contact-module}と
第\ref{sec:pf-equivariant-algebraization}節の構造的証明で得た。
これらを旧大型行列の再実行による結果とはしない。
一般の無限正規代表の収束性、固定三次式での全座標分類、
全頂点加群の低いnormal次数部分と、その乗法の完全決定は未解決である。

\subsection{再現手順と欠けている資産}
日本語LuaLaTeX原稿の組版入口は
\nolinkurl{latexmk papers/pfaffian_brill_noether_technical.tex}
である。数式・証明の現行TeXは本リポジトリの
\texttt{papers/sections/}に収録されている。
旧版の検算ドライバや証明書については、
現行Git treeに \texttt{verification/} と旧版 \texttt{source\_vNN/}
が存在しない。従ってその実行コマンドを現在の再現手順としては提示しない。
有用な検算結果・入力が回復できた場合は、旧版の一式を移すのではなく、
現行の定理を支える最小限のスクリプト、入力、証明書を選んで取り込む。
現行の\texttt{verify\_pfaffian\_vertex\_and\_tate.py}で、Schur自由分解の全記号Hilbert恒等式と、Tate再帰の六次までの有理数計算を実行した。旧大型行列の完全再実行や全証明の外部査読とは区別する。

関連先行研究の仮定と新規性、特に高次接触商と普遍共変量の既知性は
公表前に一次論文に遡って確かめる必要がある。
